\documentclass[../../../include/open-logic-section]{subfiles}


\begin{document}

\olfileid{sfr}{arith}{int}
\olsection{从 $\Nat$ 到 $\Int$}

我们有两个基本认识：
% Part: sets-functions-relations
% Chapter: arithmetization
% Section: From N to Z

\begin{enumerate}
		\item 每个整数都可以写成 $n - m$ 的形式，其中 $n, m \in\Nat$。
		\item 表达式 $n - m$ 所编码的信息，同样可以用一个有序对 $\tuple{n, m}$ 来编码。
	\end{enumerate}

我们已经知道，自然数的有序对是 $\Nat^2$ 的元素。且我们假定对 $\Nat$ 已有充分理解。于是，基于上述两个认识，这里有一个朴素的建议：\emph{让我们把整数看作自然数的有序对}。

事实上，这个建议过于朴素。显然，我们希望 $0 - 2 = 4 - 6$ 成立，但显然 $\tuple{0, 2 }\neq \tuple{4, 6}$。因此我们不能简单地说 $\Nat^2$ 就是整数的集合。

将前述问题加以推广，我们需要的是以下性质：
	$$a - b = c - d \text{ iff }a + d = c + b$$
（显然，这就是整数所\emph{应有}的性质：只需将 $b$ 和 $d$ 同时加到等式两边即可。）而要保证这种性质，一个简单的方法就是在有序对之间定义一个等价关系 $\Intequiv$，如下：
$$\tuple{ a, b } \Intequiv \tuple{c, d}\text{ iff }a + d = c + b$$
现在我们必须证明这是一个等价关系。

\begin{prop}
$\Intequiv$ 是一个等价关系。
\end{prop}

\begin{proof}
我们必须证明 $\Intequiv$ 是自反的、对称的且传递的。

\emph{自反性：} 显然 $\tuple{a, b} \Intequiv \tuple{a, b}$，因为 $a + b = b + a$。

\emph{对称性：} 设 $\tuple{a, b} \Intequiv \tuple{c, d}$，于是有 $a + d = c + b$。则 $c + b = a + d$，因此 $\tuple{c, d} \Intequiv \tuple{a, b}$。

\emph{传递性：} 设 $\tuple{a, b} \Intequiv \tuple{c, d}\Intequiv \tuple{m, n}$。于是 $a + d = c + b$ 且 $c + n = m + d$。所以 $a + d + c + n = c + b + m + d$，于是 $a + n = m + b$。因此 $\tuple{a, b } \Intequiv \tuple{m, n}$。
\end{proof}

现在我们可以用这个等价关系来取等价类：

\begin{defn}
    整数就是自然数有序对在 $\Intequiv$ 下的等价类；即，$\Int = \equivclass{\Nat^2}{\Intequiv}$。
\end{defn}

对此规定性定义，可能会有许多不同的\emph{哲学}反应。不过，在考虑这些反应之前，值得先继续处理一些技术细节。

既然已经定义了整数是什么，我们就需要定义其上的基本函数和关系。让我们用 $\equivrep{m, n}{\Intequiv}$ 来表示在 $\Intequiv$ 下包含元素 $\tuple{m, n}$ 的等价类。\footnote{注：如果使用 \olref[sfr][rel][eqv]{def:equivalenceclass} 引入的记号，同样的东西我们会写成 $\equivrep{\tuple{m,n}}{\Intequiv}$。但那样读起来稍微费劲一些。}也就是说：
    $$\equivrep{m, n}{\Intequiv} = \Setabs{\tuple{a, b} \in \Nat^2}{\tuple{a, b}\Intequiv \tuple{m, n}}$$
现在，我们给出一些定义：
    \begin{align*}
		\equivrep{a, b}{\Intequiv} + \equivrep{c, d}{\Intequiv} &= \equivrep{a + c, b + d}{\Intequiv}\\
		\equivrep{a, b}{\Intequiv} \times \equivrep{c, d}{\Intequiv} &= \equivrep{a c + b  d, a  d + b c}{\Intequiv}\\
		\equivrep{a, b}{\Intequiv} \leq \equivrep{c, d}{\Intequiv} &\text{ 当且仅当 }a + d \leq b + c
	\end{align*}
（按照惯例，我用 `$ab$' 表示 `$(a \times b)$'，仅是为了让这些定义式更易读。）现在，我们需要确保这些定义具有它们\emph{应有}的性质。详细阐明这意味着什么并进行验证相当繁琐；我们把细节放到 \olref[check]{sec}。但简而言之就是：一切正常！

还有最后一件事。我们用自然数构造了整数。但这意味着自然数\emph{本身并不是整数}。我们将在 \olref[ref]{sec} 回到这一点的哲学意义。不过，纯粹从技术角度看，我们需要某种方式能够将自然数\emph{当作}整数来处理。想法非常简单：对于每个 $n \in \Nat$，我们直接规定 $n_{\Int} = \equivrep{n, 0}{\Intequiv}$。我们需要确认这个定义是表现良好的；也就是说，对于任意 $m, n \in \Nat$，
\begin{align*}
	(m + n)_{\Int} &= m_{\Int} + n_{\Int}\\
	(m \times n)_{\Int} &= m_{\Int} \times n_{\Int}\\
	m \leq n &\liff m_{\Int} \leq n_{\Int}
\end{align*}
但这都相当简单。例如，为了证明其中第二个条件成立，我们可以直接利用自然数的性质进行如下推理：
%\begin{proof}
%	Helping myself to the behaviour of the natural numbers, evidently:
	\begin{align*}
%		(m + n)_{\Int}  = \equivrep{m + n, 0}{\Intequiv} = \equivrep{m + n, 0 + 0}{\Intequiv} = \equivrep{m, 0}{\Intequiv} + \equivrep{n, 0}{\Intequiv} = m_{\Int} + n_{\Int}\\
		(m \times n)_{\Int}  &= \equivrep{m \times n, 0}{\Intequiv} \\
		&= \equivrep{m \times n + 0 \times 0, m \times 0 + 0 \times n}{\Intequiv} \\
		&= \equivrep{m, 0}{\Intequiv} \times \equivrep{n, 0}{\Intequiv} \\
		&= m_{\Int} \times n_{\Int}
	\end{align*}
%\end{proof}
我们将其余两个条件成立的验证留作练习。

\begin{prob}
	证明对任意 $m, n \in \Nat$，有 $(m + n)_{\Int} = m_{\Int} + n_{\Int}$ 和 $m \leq n \liff m_{\Int} \leq n_{\Int}$。
\end{prob}
\end{document}